% from ega4-2-fr.tex
\oldpage[IV-2]{216}

(x) Soit $A$ un anneau local excellent. Pour tout anneau quotient intègre $B$
de $A$, $\widehat B$ est équidimensionnel. Pour que $A$ soit équidimensionnel,
il faut et il suffit que $\widehat A$ le soit.
\end{scholium}

\begin{proof}
La plupart de ces résultats ont déjà été démontrés.

(i) La première assertion résulte de (5.6.3, (i)) et (7.4.4, (ii)); la seconde,
de (7.4.4), (7.3.18), (6.12.7) et (6.12.4).

L'assertion (ii) résulte de (5.6.1), (5.6.3, (i)), (7.4.4) et (6.12.4).

(iii) La première assertion a été vue dans (5.6.4) et (7.4.4), compte tenu de
(i). La seconde a été vue dans (5.6.4) et (7.3.19, (iv)), compte tenu de (i).

L'assertion (iv) résulte de (6.12.4), (6.13.5), (6.12.9) et (6.11.8) (compte
tenu de (ii)).

La première assertion de (v) est conséquence de (7.4.6); les relations
(7.8.3.1) en découlent, compte tenu respectivement de (6.5.1), (6.5.4),
(6.5.3), (6.3.5) et (6.4.2). La dernière assertion de (v) est le cas
particulier correspondant à la propriété $(\mathrm S_n)$ pour $n=1$ (5.7.5).

L'assertion (vi) résulte de (7.7.3), et l'assertion (vii) de (7.6.1), (7.6.2)
et (7.6.3). Pour prouver (viii), notons que l'anneau $A^{(1)}$, intersection
des $A_{\mathfrak p}$ pour tous les idéaux premiers $\mathfrak p$ de hauteur
$1$ de $A$, est une $A$-algèbre finie, comme il résulte de (vi), de (7.8.2,
(i)) et de (5.11.2). Comme la clôture intégrale $A'$ de $A$ est une $A$-algèbre
finie, les idéaux premiers de hauteur $1$ de $A'$ sont exactement ceux qui sont
au-dessus des idéaux premiers de hauteur $1$ de $A$ (5.10.17, (iv)); la
conclusion résulte donc de (0, 23.2.9).

La première assertion de (ix) est conséquence de (vi) et de (5.11.4); la
seconde en est un cas particulier, si l'on observe que lorsque $\mathcal F$ est
sans torsion, $\operatorname{Ass}(\mathcal F)$ est réduit au point générique de
$\operatorname{Spec}(A)$. Enfin, pour prouver (x), il suffit d'observer que $B$
est un anneau local excellent par (ii); comme $B$ est universellement caténaire
et universellement japonais en vertu de (vi), l'anneau $B^{(1)}$ est une
$B$-algèbre finie (5.11.2), ce qui montre que $A$ est strictement formellement
caténaire par (7.2.5, $c)$) et achève la démonstration de (x).
\end{proof}

\begin{remarks}[7.8.4]
\EGAReaderTarget{ega4.c04.item.7-8-4}{IV.2.s7.8.4-fr}
\upshape
(i) Si on considère les anneaux noethériens $A$ qui vérifient seulement les
conditions (ii) et (iii) de la définition (7.8.2), on vérifie par inspection
des démonstrations précédentes que les assertions (ii), (iii), (iv), (v), (vi)
et (vii) de (7.8.3) sont encore valables, en remplaçant « excellent » par
« vérifiant » (7.8.2, (ii) et (iii)).

(ii) L'anneau local caténaire $A$ construit dans (5.6.11) vérifie les
conditions (ii) et (iii) (mais non la condition (i)) de la définition (7.8.2).
En effet, la fibre formelle en l'idéal maximal $\mathfrak n C_{\mathfrak n}$ de
$A$ est trivialement géométriquement régulière, et celle aux autres idéaux
premiers $\mathfrak p$ de $A$ coïncide avec celle de l'anneau $E$ en ces idéaux
premiers; or $E$ est une $V$-algèbre de type fini et $V$ est un anneau de
valuation discrète, qui est par suite un anneau excellent si l'on prend le
corps $k_0$ de caractéristique $0$ (7.8.3, (iii)); par suite $E$ est un anneau
excellent (7.8.3, (ii)) et cela prouve donc que $A$ vérifie la condition (7.8.2,
(ii)). D'autre part, le complémentaire du point fermé dans
$\operatorname{Spec}(A)$ est isomorphe à un ouvert de $\operatorname{Spec}(E)$,
dont un ouvert affine est le spectre d'un anneau

